\chapter{Zusammenfassung und Ausblick}
\label{sec:zusammenfassung}

In der vorliegenden Bachelorarbeit wird ein robustes Zugangssystem für Katzen entwickelt, dass das eindringen von Beutetieren (z. B. Mäuse oder Vögel) in den Wohnraum verhindert. Durch die Kombination von mehrschitigen Lernen (eng. Deep Learning) und zeitlicher Aggregation wurde ein System geschaffen, das die einzelbildbasierte Objekterkennung auf den Testdaten deutlich übertrifft.

\section{Zusammenfassung der Erkenntnisse}
\label{subsec:zusammenfassung_erkenntnisse}

Die theoretischen Annahmen konnten durch die praktische Evaluierung erfolgreich validiert werden. In diesem Anwendungsfall ist der klassische monolithische Objektdetektor deutlich ungenauer als die in der Arbeit vorgestellte Verarbeitungskette, die aus einem Objektdetektor, einer Bildsegmentierung und einer zeitlichen Aggregation besteht.

Zunächst detektiert ein YOLO-Detektor das Gesicht der Katze. Um zu verhindern, dass die eventuelle Beute im Maul der Katze okkludiert wird, wird der Ausschnitt des Gesichts um rund eine Gesichtshöhe nach unten erweitert. Im zweiten Schritt analysiert ein YOLO-Segmentierungsmodell, ob sich Beute im Maul der Katze befindet. Um Fehlklassifikationen auszugleichen, muss eine Katze beweisen, dass sie sauber ist. Erst dann darf sie in den Wohnraum eintreten. Dazu wird die Katze über mehrere Sekunden hinweg mithilfe von ByteTrack verfolgt. Der gleitende Durchschnitt der letzten 30 Einzelbilder (was bei 8 Bildern pro Sekunde 3,75 Sekunden entspricht) entscheidet dann, ob die Katze ins Haus darf oder nicht.

Die Evaluierung mit KI-generierten Videodaten zeigte den Erfolg des gewählten Architektur-Ansatzes. Das System gab die Klappe in keinem Einzelbild mit Beute fälschlicherweise frei und sperrte keine saubere Katze aus. Der gewählte Schwellenwert von 0,3 sperrt die Klappe bereits bei geringen Anzeichen von Beute. Saubere Katzen warten dafür im Mittel 2,0 Sekunden, bis die Klappe freigegeben wird. Das System erweist sich somit als sehr praxistauglich und läuft mit durchschnittlich 8 Bildern pro Sekunde ressourcenschonend auf lokalen Endgeräten.

Somit konnte die zentrale Forschungsfrage bejaht werden. Eine Kombination aus Objekterkennung, Objektverfolgung und zeitlicher Glättung ist deutlich robuster als herkömmliche Einzelbildauswertungen.

\section{Ausblick und Weiterentwicklungsmöglichkeiten}
\label{subsec:ausblick}

Obwohl das System in seiner aktuellen Form bereits zuverlässig Beutetiere aussperrt, zeigt die Praxis, dass auch dieser Ansatz Schwächen und Limitierungen hat. Das eröffnet die Möglichkeit, konkrete Ansatzpunkte für zukünftige Verbesserungen zu identifizieren.

\subsection{Kontinuierliches Sammeln unsicherer Beispiele}
Die vielversprechendste Methode, um die Systemleistung zu steigern, ist die Erweiterung des Datensatzes für Katzen mit Beute. Um dieses Ziel zu automatisieren, verfügt die Applikation über eine Funktionalität, die Bilder speichert, bei denen die Bildsegmentierung unsicher war. Dadurch können zukünftige Arbeiten diesen stetig wachsenden Datensatz nutzen, um das YOLO-Segmentierungsmodell iterativ nachzutrainieren. Sobald das Segmentierungsmodell aus diesen Fehlern lernen kann, kann die Fehlalarmrate des gesamten Systems sinken. Dieses Vorgehen entspricht dem Prinzip des aktiven Lernens, bei dem gezielt jene Beispiele ausgewählt werden, bei denen das Modell am unsichersten ist \cite{settles2009active}.

\subsection{Verlust der Tiefeninformationen}
Eine weitere Einschränkung ist der Verlust von 3D-Informationen im Bild. Dadurch kann es passieren, dass Objekte im Hintergrund oder sogar der Körper der Katze fälschlicherweise als \glqq im Maul der Katze\grqq\ interpretiert werden. Eine mögliche Erweiterung bestünde darin, diese 3D-Informationen mithilfe von Stereo-Kameras oder Time-of-Flight-Sensoren zurückzugewinnen. Dadurch könnten Fehlalarme herausgefiltert werden, bei denen sich das vermeintliche Beutetier weit hinter dem Maul der Katze befindet.

\subsection{Hardware-Beschleunigung}
Derzeit wird das System ausschließlich auf der CPU ausgeführt. Ein Wechsel auf eine dedizierte Hardware (z. B. Google Coral) würde die Anzahl der pro Sekunde verarbeiteten Bilder deutlich erhöhen. Dadurch können wir in der gleichen Zeit mehr Daten der Katze sammeln, was die zeitliche Glättung verbessern würde. Das Flackern des Systems würde dadurch noch weiter eingeschränkt.

\vspace{0.5cm}

Abschließend lässt sich zusammenfassen, dass die in dieser Arbeit entwickelte Verarbeitungskette ein kompaktes und robustes Fundament bildet. Sie löst das Problem des unerwünschten Beuteeintrags durch Katzen zuverlässig und ermöglicht zugleich spannende Erweiterungen und Verbesserungen.
